\documentclass[10pt,a4paper]{amsart}
\usepackage[margin=19mm]{geometry}
\usepackage{amsmath,amssymb,mathtools}
\usepackage[scr=rsfs,cal=euler]{mathalfa}
\usepackage{xfrac}
\usepackage[unicode,hidelinks,bookmarks=false]{hyperref}
\newcommand{\ie}{i.e.\ }
\newcommand{\cf}{cf.\ }
\newcommand{\resp}{respectively\ }
\newcommand{\Subsection}{\subsection}
\providecommand{\nobreakdash}{\nobreak\hbox{-}}
\setlength{\emergencystretch}{2em}
\multlinegap0pt
\usepackage{amssymb}
\newtheorem{proposition}{Proposition}
\newcommand{\e}{\mathrm e}
\title{A counterexample to the zero-mass conjecture}
\author{Long Li, Mingchen Xia}
\date{Source: arXiv:2607.26549v1 (CC BY 4.0)}
\begin{document}
\maketitle

\noindent\emph{Source and licence.} A counterexample to the zero-mass conjecture. Version arXiv:2607.26549v1 (CC BY 4.0), distributed by arXiv under the Creative Commons Attribution 4.0 International licence, http://creativecommons.org/licenses/by/4.0/, as recorded in the arXiv licence metadata for this exact version and shown on the abstract page https://arxiv.org/abs/2607.26549v1. Full licence notice: \texttt{licenses/arxiv-2607.26549-CC-BY-4.0.txt}.\par\smallskip

\noindent\textit{Editorial scope.} Counted unit: the article's proposition prop:limit (source0.tex lines 426--433). The article's complete original proof of it is delivered with it (source0.tex lines 435--473, one \texttt{proof} environment of 39 lines). No mathematics is written, repaired or re-derived: the statement and the proof are the article's own lines, in the article's own notation, and no retained line is substituted or re-flowed.  Retained uncounted as the source-local context the proof needs: lines 336--338; lines 367--369; lines 371--373; lines 414--416; lines 422--424.  Nothing was omitted from any retained span, so the package carries no omission record.  No retained line contains a citation, so the wrapper carries no bibliography and no bibliography entry.  No retained line contains a figure environment, an image-inclusion command or drawing code, so no image is delivered and the package declares no asset dependency.  Editorial formatting only, in addition: an AMS \texttt{amsart} preamble that restates the article's own declarations and macros used by the retained lines, namely the environments \texttt{proposition} on separate counters, together with the standard packages those lines need.  The retained environments print in this document as Proposition 1 in order of appearance, whereas the article prints the same items with the numbering of its own full document.  The complete native source of the article as distributed in the arXiv e-print for this exact version is bundled unchanged as \texttt{sources/206278/source0.tex}.
\begin{equation}\label{eq:q-lower}
  q_j(\zeta)\ge 2^j\log\|\zeta\|_\infty-(\log 2)(2^j-2^{-j}).
\end{equation}

\begin{equation}\label{eq:q-visible}
  -2^{-j}\log 2 \le q_{j+1}-q_j \le 0 \quad\textup{on }\{q_j\ge-4^j\},
\end{equation}

\begin{equation}\label{eq:q-global}
  q_{j+1}\le\max\{q_j,-4^j\} \quad\textup{on }\mathbb D^2.
\end{equation}

\begin{equation}\label{eq:Vj}
  V_j=\max\{q_j,-4^j\}
\end{equation}

\begin{equation}\label{eq:u}
  u=\inf_{j\ge1}V_j.
\end{equation}

\begin{proposition}\label{prop:limit}
Each $V_j$ is a bounded negative plurisubharmonic function, and the sequence decreases pointwise to $u$.
The limit is negative and plurisubharmonic on $\mathbb D^2$, continuous on $\mathbb D^2\setminus\{0\}$, and
\[
  u^{-1}(-\infty)=\{0\}.
\]
Moreover, $q_j\to u$ locally uniformly on $\mathbb D^2\setminus\{0\}$, and $\e^u$ is continuous on $\mathbb D^2$.
\end{proposition}

\begin{proof}
Each $V_j$ is bounded and plurisubharmonic. Since $H_j(\mathbb D^2)\subset\mathbb D^2$, both branches in \eqref{eq:Vj} are negative, so $V_j<0$. Moreover, \eqref{eq:q-global} gives $q_{j+1}\le V_j$, while
\[
  -4^{j+1}\le-4^j\le V_j.
\]
Taking the maximum of these two inequalities proves $V_{j+1}\le V_j$. Thus $V_j$ decreases pointwise to the function $u$ defined in \eqref{eq:u}.

If $\zeta\in K_j$, then \eqref{eq:q-lower} gives
\[
  2^j\log\|\zeta\|_\infty \le-4^j+(\log 2)\left(2^j-2^{-j}\right) \le-4^j+2^j\log 2,
\]
and therefore
\begin{equation}\label{eq:core}
  \|\zeta\|_\infty
  \le 2\e^{-2^j}.
\end{equation}

Moreover, $K_j=\{\zeta\in\overline{\mathbb D^2}:\e^{q_j(\zeta)}\le\e^{-4^j}\}$ is closed in $\overline{\mathbb D^2}$,
because $\e^{q_j}$ is continuous there. Thus the cores lie in closed max-norm balls whose radii tend to zero. Since $2\e^{-2}<1$, every $K_j$ is compactly contained in $\mathbb D^2$.

Let $L\Subset \mathbb D^2\setminus\{0\}$. By \eqref{eq:core}, there is an $N$ such that $L\cap K_k=\varnothing$ for every $k\ge N$. Hence \eqref{eq:q-visible} applies at each later stage, and for $\ell>k\ge N$,
\[
  0\le q_k-q_\ell
  \le\sum_{s=k}^{\ell-1}2^{-s}\log 2
  \le2^{1-k}\log 2
  \qquad\text{on }L.
\]
Thus $q_j$ is uniformly Cauchy on $L$. The cutoff in \eqref{eq:Vj} is also inactive on $L$ for $j\ge N$, so $V_j=q_j$ there. Therefore $V_j$ and $q_j$ have the same locally uniform limit on the punctured domain.

A decreasing limit of plurisubharmonic functions is plurisubharmonic unless it is identically $-\infty$. The punctured limit is finite, so that alternative is excluded, and $u\le V_1<0$. Finally,
\[
  V_j(0)=-4^j\longrightarrow-\infty.
\]
The locally uniform limit of the continuous functions $q_j$ is continuous away from $0$. To prove continuity at the pole, fix $A>0$ and choose $j$ with $4^j>A$. Since $\e^{q_j}$ is continuous and vanishes only at $0$, the core $K_j$ contains a neighborhood of $0$. On this neighborhood,
\[
  u\le V_j=-4^j<-A.
\]
Thus $u(\zeta)\to-\infty$ as $\zeta\to0$, so $\e^u$ is continuous at $0$ as well.
\end{proof}

\end{document}
